% <= 400 words (guidelines). Numbers: n04, n05 (H3), n06.
Do the qualities journalists call newsworthy predict what readers click, and does public importance align with engagement? We scored all 22,771 candidate articles of the MIND-small news recommendation dataset on seven news values (conflict, magnitude, prominence, surprise, utility, entertainment and public importance) with an open-weight large language model, Qwen2.5-7B-Instruct, from headline and abstract. Against two independent human annotators, six values passed a preregistered gate (inter-annotator $\kappa$ 0.61 to 0.81; LLM-human $\rho$ 0.46 to 0.78); entertainment failed on annotator agreement. We then related the scores to clicks within 230,117 impressions, comparing only items shown together, so that exposure differences cancel. We show that the usual user-level bootstrap is badly anti-conservative for article attributes, since a random per-article score appears significant, and use a two-way users-by-articles bootstrap and article permutation tests instead. The preregistered hypothesis that prominence and surprise predict clicks was not supported. Public importance was unrelated to clicks (within-impression AUC 0.505, 95\% CI [0.467, 0.548]; article-level $\rho$ with click rate $-$0.008). Practical utility predicted fewer clicks (AUC 0.434 [0.392, 0.456]) and conflict more (0.535 [0.510, 0.572]), both replicating on a second data split. All effects were small next to past popularity (AUC 0.587). Engagement is therefore not a proxy for importance, and item-level analyses of click logs need article-aware inference.
